% ---------------------------------------------------------------------
%  04 · Composants — les « signatures » du design system
%
%    \pastille{2}                    numéro dans un rond herbe
%    \etiquette[1]{Le plat}          pilule neutre, numéro optionnel
%    \badge[agrume]{25 min}          pilule colorée : agrume | herbe | sauge | contour | encre
%    \bandeau[centre]{gauche}{droite} bandeau daté : filet épais, capitales, filet fin
%    \pointilles                     points de conduite (remplit la ligne)
%    lignes + \ligne[n]{libellé}{valeur}   « 2  Le verre ........ Vermentino »
%    carte[Libellé] · note[Libellé] · attention[Libellé]   encadrés arrondis
% ---------------------------------------------------------------------

% Pastille numérotée ---------------------------------------------------
\newcommand\pastille[2][1.5em]{%
  \tikz[baseline=(p.base)]
    \node (p) [circle, fill=action, text=surAction, inner sep=0pt,
               minimum size=#1,
               font=\policeChiffres\bfseries\fontsize{0.62em}{0.62em}\selectfont]
               {#2};%
}

% Pilules --------------------------------------------------------------
\tikzset{
  pilule/.style = {rounded corners=0.8em, minimum height=1.6em,
                   inner xsep=0.7em, inner ysep=0pt,
                   text height=1.6ex, text depth=0.25ex,
                   font=\policeMoyenne\bfseries\small},
  pilule agrume/.style  = {pilule, fill=accent,       text=surAccent},
  pilule herbe/.style   = {pilule, fill=teinteAction, text=action},
  pilule sauge/.style   = {pilule, fill=surface,      text=texte2},
  pilule contour/.style = {pilule, draw=texte, line width=0.7pt, text=texte},
  pilule encre/.style   = {pilule, fill=texte,        text=fond},
}

\newcommand\badge[2][agrume]{%
  \tikz[baseline=(b.base)]\node (b) [pilule #1] {#2};%
}

\NewDocumentCommand\etiquette{o m}{%
  \tikz[baseline=(e.base)]\node (e) [pilule, fill=surface, text=texte]
    {\IfValueT{#1}{\textcolor{action}{#1}\hspace{0.45em}}#2};%
}

% Bandeau daté ---------------------------------------------------------
\NewDocumentCommand\bandeau{O{} m m}{%
  \vbox{%
    {\color{filetFort}\hrule height 1.1pt}%
    \vskip 6pt
    \hbox to \hsize{%
      \strut\textlabel{#2}\hfill
      \if\relax\detokenize{#1}\relax\else\textlabel{#1}\hfill\fi
      \textlabel{#3}}%
    \vskip 4pt
    {\color{filet}\hrule height 0.5pt}%
  }%
}

% Points de conduite ---------------------------------------------------
\newcommand\pointilles{%
  \leaders\hbox to 0.42em{\hss\textcolor{texte2!70}{.}\hss}\hfill
}

% Lignes numérotées à pointillés ---------------------------------------
\newenvironment{lignes}{%
  \par\addvspace{0.8em}%
  {\color{filet}\hrule height 0.5pt}%
}{%
  \par\addvspace{0.8em}%
}

\NewDocumentCommand\ligne{o m m}{%
  \par\vskip 5pt
  \hbox to \linewidth{%
    \strut
    \IfValueT{#1}{\pastille{#1}\hspace{0.7em}}%
    {\policeMoyenne\color{texte2}#2}%
    \hspace{0.6em}\pointilles\hspace{0.6em}%
    {\policeIntertitre\large #3}}%
  \vskip 5pt
  {\color{filet}\hrule height 0.5pt}%
}

% Encadrés -------------------------------------------------------------
\tcbset{
  marche/.style = {
    enhanced, breakable, frame hidden,
    arc=3mm, boxrule=0pt,
    colback=surface, coltext=texte,
    left=5mm, right=5mm, top=4.5mm, bottom=4.5mm,
    before skip=1em, after skip=1em,
    parbox=false,
  },
  libelle/.style = {
    before upper={%
      \if\relax\detokenize{#1}\relax\else
        \textlabel{#1}\par\vspace{0.3em}%
      \fi},
  },
}

\newtcolorbox{carte}[1][]{marche, libelle={#1}}

\newtcolorbox{note}[1][Note]{
  marche, colback=teinteAction,
  before upper={\textlabel[action]{#1}\par\vspace{0.3em}},
}

\newtcolorbox{attention}[1][Attention]{
  marche, colback=teinteAccent,
  before upper={\textlabel{#1}\par\vspace{0.3em}},
}

% Trou à combler : visible dans le PDF tant que le contenu manque
\newcommand\acompleter[1]{%
  \begin{attention}[À compléter]#1\end{attention}%
}
